\subsection{An odd-order diamond at level \texorpdfstring{$135$}{135}}
\begin{proposition}\label{prop:decic135}
No elliptic curve over a decic number field has a point of order $135$.
\end{proposition}
\begin{proof}
Put $X=X_1(135)$, $J=\operatorname{Jac}(X)$, and
$B=(\langle16\rangle-1)J$. The complete winding computation has
weighted dimension $529$, equal to the genus of $X$. Every nontrivial
character component is generated by its winding element. Thus
$B(\Q)$ is finite.

The character components contributing to $B$ have orders
$9,3,36,12,18,6$, and respective dimensions
$16,12,16,12,16,12$ over their character fields. Their total rational
dimension is $480$. For each component we compute
\[
 \operatorname{Norm}\det\bigl((1+31\chi(31))I-T_{31}\bigr).
\]
All six norms are odd. Hence $\#B(\mathbb F_{31})$ is odd, and
$B(\Q)[2^\infty]=0$. The complete rational matrices and their norms
are replayed independently in Sage.

The index is $7776$, and the global criterion at $2$ has $k_2=6$.
Its required inequality is
\[
 60<\frac{325\cdot7776}{32768}=\frac{78975}{1024}.
\]
The minimum CM degree is $72$. A hypothetical noncuspidal decic
trace divisor must therefore be invariant under $\langle16\rangle$.
This diamond has projective order $9$ and acts freely on the
non-CM locus. Its orbit lengths cannot sum to $10$.
\end{proof}

\subsection{Cuspidal reduction and relative quintic divisors}
\begin{lemma}\label{lem:decic-cuspidal-reduction}
Let $X/F$ be a smooth projective curve with a rational cusp $c_0$
and rank-zero Jacobian. Suppose that its cuspidal divisor classes
are torsion. Let $v$ be an unramified prime of good reduction,
and suppose the cusps split over an extension unramified at $v$.
Let $D$ be an effective $F$-rational divisor of degree $d$, with
noncuspidal support, whose reduction is entirely cuspidal.
\begin{enumerate}
\item If the residue characteristic is odd, then
$\ell(D)\geq2$.
\item If the residue characteristic is $2$, then
$\ell(2D)\geq2$. More precisely, if every effective cuspidal residue
divisor $C$ of degree $d$ has $\ell(2C)=1$, no such $D$ exists.
\end{enumerate}
\end{lemma}
\begin{proof}
After an unramified extension lift the reduction to an effective
cuspidal divisor $C$ of degree $d$. The classes of $D-dc_0$ and
$C-dc_0$ are torsion, by rank zero and Manin--Drinfeld, respectively.
Thus $[D-C]$ is torsion in the kernel of reduction.
At an unramified odd prime this kernel is torsion free. Hence
$D\sim C$, and the two distinct effective divisors give
$\ell(D)\geq2$, also over $F$ by cohomology base change.

At an unramified dyadic prime the torsion reduction kernel has
exponent two: multiplication by two carries the first formal-group
filtration into the second, on which the formal logarithm is
injective. Hence $2D\sim2C$, so $\ell(2D)\geq2$.
Specialization gives $\ell(2\overline C)\geq2$ as well.
\end{proof}

\begin{proposition}\label{prop:decic321}
No elliptic curve over a decic number field contains
$C_3\times C_{21}$.
\end{proposition}
\begin{proof}
Put $F=\Q(\zeta_3)$ and $X=X_1(3,21)$. The Weil pairing gives a
noncuspidal effective $F$-rational divisor $D$ of degree $5$.
The curve has genus $25$, and its Jacobian has rank zero over $F$
by \cite[Theorem~4.1]{DS}. Its index is $576$, so the Abramovich
bound gives $\operatorname{gon}_F X>5$.

Over $\mathbb F_4$ the cusp signature is $1^{24},3^8$,
and the noncuspidal closed-point counts in degrees one through five
are
\[
 (0,0,0,144,216).
\]
The Frobenius polynomial of $J_H(63)$, with $H=\langle22\rangle$,
independently gives these counts after quadratic constant-field
extension. There are
\[
 \binom{28}{5}+8\binom{25}{2}=100680
\]
effective cuspidal degree-$5$ divisors. For all of them the exact
local differential matrices have combined rank $10$, giving
$\ell(2C)=1$. There are $128$ local matrices, including every
admissible nonreduced multiplicity. The native calculation checks
each local dimension by ordinary Riemann--Roch. An independent
exhaustive $\mathbb F_4$ replay checks all $100680$ combined ranks.
Lemma~\ref{lem:decic-cuspidal-reduction} excludes entirely cuspidal
reduction.

Every rational residue cusp is defined over $F$; this follows
from the complete primitive Tate-cusp permutation under
$\operatorname{Gal}(\Q(\zeta_{21})/F)$.
For $A_{13}=T_{13}-1-13\langle13\rangle$, the split-prime
Eichler--Shimura relation and odd torsion injectivity give
$A_{13}J(F)=0$. The cuspidal boundary formula gives $A_{13}c=0$
for every rational cusp. Thus a quartic-plus-cusp reduction requires
$A_{13}x$ principal at its quartic support.

The pairing-preserving level group
$(\operatorname{SL}_2(\mathbb F_3)\times\mathbb F_7^\times)
/\langle(-I,-1)\rangle$ has order $72$ and commutes with $A_{13}$.
The $144$ quartic places form two orbits, and the $216$ quintic places
form three orbits. At all five representatives, $A_{13}x$ is
nonprincipal. The quartic tests are the completed degree-independent
calculations preserved from the earlier work. The new quintic
calculation checks every orbit and compares the first Hecke divisor
with the generic isogeny construction; its degree is $70$.
These tests exclude every possible reduction of $D$.
\end{proof}

\begin{proposition}\label{prop:decic324424}
No elliptic curve over a decic number field contains either
$C_3\times C_{24}$ or $C_4\times C_{24}$.
\end{proposition}
\begin{proof}
For $m=3,4$ put $F=\Q(\zeta_m)$ and $X=X_1(m,24)$.
A hypothetical subgroup gives a noncuspidal effective divisor
$D$ of degree $5$ over $F$. Both Jacobians have rank zero over $F$
by \cite[Theorem~4.1]{DS}. Their indices are $576$ and $768$,
respectively, so both have $F$-gonality greater than $5$.
The cusps split in $\Q(\zeta_{24})$.

For $m=3$ we reduce at $7$, and for $m=4$ at $5$.
These primes split in $F$. Lemma~\ref{lem:decic-cuspidal-reduction}
excludes entirely cuspidal reduction. The complete local counts are
\[
\begin{array}{c|c|c|rrrrr}
m&q&\text{cusp signature}&\multicolumn{5}{c}{\text{noncuspidal degrees }1,\ldots,5}\\\hline
3&7&1^{24},2^{12}&0&0&48&456&3456\\
4&5&1^{16},2^{24}&0&0&0&96&704
\end{array}
\]
Independent Frobenius polynomials of $J_H(72,\langle25\rangle)$
and $J_H(96,\langle25\rangle)$ verify the counts.
Every rational residue cusp is globally rational over $F$, as
checked by the full Tate-cusp Galois action.

Suppose first that $m=4$. The reduction is a quartic place and a
rational cusp, or one quintic place. The $64$ pairing-preserving
level automorphisms partition these supports into four and eleven
orbits, respectively. At every representative the exact divisor
$A_{13}x$ is nonprincipal. Since $A_{13}$ kills $J(F)$ and every
rational cusp, this is impossible.

Now let $m=3$. The $48$ level automorphisms partition the cubic,
quartic and quintic supports into $1$, $13$ and $72$ orbits.
All quartic and quintic representatives have nonprincipal
$A_{13}$-image. This excludes the quartic-plus-rational-cusp and
quintic branches.

A cubic support can be accompanied by any effective cuspidal
divisor of degree two. To cover all such complements, use
$A_{31}=T_{31}-1-31\langle31\rangle$.
For a geometric cusp the boundary relation is
\[
 T_q c=\sigma_q c+q\langle q\rangle\sigma_q^{-1}c,
\]
where $\sigma_q$ is cyclotomic Frobenius.
Since $31\equiv7\pmod{24}$, $\sigma_{31}$ fixes every closed
cuspidal divisor over $\mathbb F_7$. Hence $A_{31}$ kills every
local cusp as a divisor, including the quadratic cusps.
The prime $31$ splits in $F$, so $A_{31}J(F)=0$.
At the single cubic orbit representative, $A_{31}x$ is
nonprincipal. This excludes the cubic branch.

The exported orbit partitions cover all actual closed places.
Independent replay checks their exact coverage across the parallel
jobs, all point counts, and every recorded principality decision.
The Hecke algorithm verifies the cyclic-kernel count and all
residue-degree multiplicities.
\end{proof}

\subsection{Quadratic descent at level \texorpdfstring{$155$}{155}}
\begin{proposition}\label{prop:decic155}
No elliptic curve over a decic number field has a point of order $155$.
\end{proposition}
\begin{proof}
The complete winding computation at level $155$ has weighted
dimension $841$, equal to the genus. The only components not
fully generated by their winding elements have character orders
$1$, $12$ and $3$. The diamond $\delta=\langle16\rangle$
has order $5$, so it acts trivially on these components.
Thus $(\delta-1)J_1(155)(\Q)$ is finite.
The global criterion at $3$ applies, since
\[
 80<\frac{325\cdot11520}{32768}=\frac{14625}{128}.
\]
The minimum CM degree is $120$.

Let $x$ be a proposed noncuspidal point. Its residue field $K$
has degree $10$, since degrees $1,2,5$ are excluded by the known
classifications. Diamond invariance gives an order-five
$\Q$-automorphism of $K$. Its fixed field $F$ is quadratic.
Since $16\equiv1\pmod5$, the image on $X_1(5)$ is defined over $F$.
The fine moduli interpretation descends the elliptic curve together
with its point of order $5$ to $F$. The cyclic subgroup of order $31$
is preserved as well, and therefore gives an $F$-rational $31$-isogeny.

Put
\[
 Y=X_1(31)/\langle2\rangle
   =X_{\{\pm1,\pm2,\pm4,\pm8,\pm16\}}(31).
\]
Since $16=2^4\pmod{31}$, the image $y$ on $Y$ is defined over $F$.
Its trace gives a noncuspidal effective $\Q$-rational divisor $D$
of degree two. The curve $Y$ has genus $6$ and a rational cusp.
Its Jacobian has rank zero, as a quotient of $J_1(31)$; the complete
winding computation at level $31$ has dimension $26$ and is full
on every character component. By \cite[Proposition~3.6]{Orlic},
\[
 \operatorname{gon}_{\Q}Y\geq
 \operatorname{gon}_{\mathbb F_2}Y\geq5.
\]

We show that $y$ has cuspidal reduction at every prime $v\mid2$
of $F$. An elliptic curve with an $F_v$-rational point of order $5$
cannot have additive reduction: the smooth identity component is
an additive group in characteristic two, and the geometric component
group has order at most four. The formal group has no $5$-torsion.
Thus the curve has good or multiplicative reduction.
If $y$ had noncuspidal reduction, its $j$-invariant would be integral,
so the reduction would be good.

The residue field then has order $q=2$ or $4$. Reduction preserves
the point of order $5$, and the $31$-isogeny supplies a
Frobenius-invariant line in $E[31]$. Hasse's bound forces
$\#E(\mathbb F_q)=5$, giving Frobenius trace $-2$ or $0$.
The corresponding characteristic-polynomial discriminants are
$-4$ and $-16$. Both are nonsquares modulo $31$, so no invariant
line exists. This contradiction proves the claimed cuspidal reduction.

The curve $Y$ has good reduction at $2$, and its cusps split in
$\Q(\zeta_{31})$, which is unramified there.
Lemma~\ref{lem:decic-cuspidal-reduction} now gives
$\ell(2D)\geq2$. Since $2D$ is a rational divisor of degree four,
this contradicts $\operatorname{gon}_{\Q}Y\geq5$.
\end{proof}
